%
%$Log: taguide.tex,v $
%Revision 1.6  2002/10/21 15:26:44  mjm
%to ohio, revised
%
%Revision 1.5  2001/08/15 18:07:03  mjm
%APPM preprint 466
%
%Revision 1.4  1999/08/09 19:55:27  mjm
%beta release
%
%Revision 1.3  1999/04/06 19:29:00  mjm
%Matts comments
%
%Revision 1.2  1999/04/06 19:22:02  mjm
%Ellens comments
%
%


%%%%%%%%%%%%%%%%%%%%%%%%%%%%%%%%%%%%%%%%%%%%%%%%%%%%%%%%%%%%%%%%%%%%%%%%%
%\handout{Teaching Assistant's Guide to Good Problems}
\handout{TA's Guide to Good Problems}

This guide is for the teaching assistant, who runs the recitations.
It lists those tasks that the teaching assistant must perform, indicates
when they must be performed, and estimates the time required.
If the course is taught without recitations, then these tasks
fall on the instructor.

The first time you use this method, there are extra
tasks, which we list first.
\begin{enumerate}
	\item Read and understand the material in this packet. This
should be performed in the week before the term starts.
Time required: 2 hours.
	\item Attend training. The instructor should provide some training
during the week before the term starts. Time required: 1 hour.
	\item Deal with unanticipated problems. Time required: 2 hours.
\end{enumerate}

Each time the method is used there are the following tasks.
\begin{enumerate}
	\item Prepare for the new skill to be taught that week. It may also
be useful to review the old skills. This is part of the weekly preparation
time. Time required: 5 minutes per week.
	\item Present difficult skills in the recitation. Sometimes
students will have particular difficulty on some specific skill. The
teaching assistant will need to re-explain this skill and perhaps
demonstrate it on a relevant problem. It is important to note that the
teaching assistant is not expected to be the main teacher of these
skills. The handouts should provide all the information the students
need. Time required: 5 minutes recitation time per week, on average.
	\item Grade the good problems. It seems fair to
grade two less regular homework problems to compensate for the good
problems. After some initial trauma, it is expected that the good problems
will be more pleasant to grade than ordinary problems. There should be no
net gain in grading time as a result of this.
	\item Provide feedback on the handouts. Any deficiencies in the
method or the handouts should be noted. At the end of the term these
comments should be collected by the instructor and the materials
corrected. Time required: 1 hour. 
\end{enumerate}
